% ---------------------------------------------------------------------
% Experimental Setup (~2 pages). Counts in tab:datasets are from the
% official corpus papers; verify against your local manifests. [DUMMY]
% markers flag values to re-check or replace.
% ---------------------------------------------------------------------

\subsection{Datasets and Protocol}
\label{ssec:datasets}

Table~\ref{tab:datasets} summarizes the seven corpora in our benchmark. We follow a
strict \emph{train-once, evaluate-everywhere} protocol: all trainable systems are
trained exclusively on the ASVspoof~2019~LA training partition (plus our augmentation
curriculum, which uses no additional speech data) and model selection uses only its
development partition. The 2019~LA evaluation set serves as the in-domain test; the six
remaining corpora are evaluated \emph{zero-shot}---their synthesis algorithms,
speakers, languages, recording conditions, and channels are never seen in training.
This protocol isolates generalization, the property that in-domain leaderboards fail to
measure~\cite{muller2022inthewild,muller2024harder}.

% [DUMMY] — verify utterance counts against local manifests before submission
\begin{table}[t]
\centering
\caption{Evaluation corpora grouped by protocol role and the failure mode (G1--G3)
  each one probes.}
\label{tab:datasets}
\setlength{\tabcolsep}{3pt}
\footnotesize
\resizebox{\columnwidth}{!}{%
\begin{tabular}{lllcr}
\toprule
Dataset & Role & Probes & Utts. (spoof/total) & Lang. \\
\midrule
ASVspoof 2019 LA~\cite{wang2020asvspoof2019} & Train/dev/eval & parity & 63,882/121,461 & EN \\
ASVspoof 2021 LA~\cite{liu2023asvspoof2021}  & Zero-shot & G2 & 163,114/181,566 & EN \\
ASVspoof 2021 DF~\cite{liu2023asvspoof2021}  & Zero-shot & G1+G2 & 589,212/611,829 & EN \\
ASVspoof 5~\cite{wang2025asvspoof5journal}   & Zero-shot & G1 (2024 attacks) & 542,466/680,774 & EN \\
In-the-Wild~\cite{muller2022inthewild}       & Zero-shot & G1+G2 (found) & 11,816/31,779 & EN \\
WaveFake~\cite{frank2021wavefake}            & Zero-shot & G1 (vocoders) & 117,985/134,097 & EN, JA \\
Codecfake~\cite{xie2025codecfake}            & Zero-shot & G1 (codec TTS) & 740,747/785,591 & EN, ZH \\
MLAAD v5~\cite{muller2024mlaad}              & Zero-shot & X-lingual & 163,900/$\sim$328k & 38 \\
\bottomrule
\end{tabular}%
}
\end{table}

For MLAAD we pair spoofed utterances with bona-fide speech from M-AILABS as prescribed
by the corpus authors~\cite{muller2024mlaad}, and report both pooled EER and
per-language-family breakdowns. For ASVspoof~5 we use the official evaluation
partition, including its adversarially optimized attacks; for Codecfake we evaluate on
the codec-resynthesis test condition. Robustness sweeps (Section~\ref{ssec:robustness})
re-encode the ASVspoof~2019 evaluation set through controlled codec, noise, and
reverberation conditions so that channel effects can be measured against a fixed attack
distribution.

\subsection{Baselines}
\label{ssec:baselines}

We compare AuralGuard against eight systems spanning the four countermeasure
generations of Section~\ref{ssec:related-cm}, all retrained under our protocol unless
noted:

\begin{itemize}
    \item \textbf{B1 --- LFCC-LCNN}~\cite{lavrentyeva2019stc}: hand-crafted LFCC
        features with a light CNN (generation 1--2 reference).
    \item \textbf{B2 --- RawNet2}~\cite{tak2021rawnet2}: end-to-end raw-waveform
        SincConv network.
    \item \textbf{B3 --- AASIST}~\cite{jung2022aasist}: spectro-temporal graph
        attention on a RawNet2-style encoder.
    \item \textbf{B4 --- W2V2+AASIST}~\cite{tak2022automatic}: wav2vec~2.0 XLS-R
        front-end with AASIST back-end and RawBoost, the reference SSL recipe.
    \item \textbf{B5 --- WavLM+AASIST+OCS}: WavLM-Large front-end (last layer), AASIST
        back-end, OC-Softmax loss---the strongest single-view configuration of our own
        stack, and the direct ablation anchor for the multi-view contributions.
    \item \textbf{B6 --- XLSR-Mamba}~\cite{xiao2025xlsrmamba}: dual-column
        bidirectional state-space detector; published state of the art on
        ASVspoof~2021~LA and In-the-Wild.
    \item \textbf{B7 --- SLIM}~\cite{zhu2024slim}: style--linguistics mismatch model
        (numbers from the original paper where retraining is not possible; marked
        $\dagger$ in tables).
    \item \textbf{B8 --- ALLM zero-shot}~\cite{gong2025allm4add,chu2024qwen2audio}:
        Qwen2-Audio-7B-Instruct prompted with the binary detection question of
        ALLM4ADD (``Does this audio sound like real human speech or AI-generated
        speech?''), scored via the token probability of the answer; no task-specific
        training. This quantifies how far general-purpose audio foundation models are
        from specialized countermeasures.
\end{itemize}

\subsection{Implementation Details}
\label{ssec:implementation}

% [DUMMY] — confirm final hyperparameters from config/ before submission
\begin{table}[t]
\centering
\caption{Training hyperparameters (selected on the ASVspoof~2019~LA dev set).}
\label{tab:hyper}
\footnotesize
\setlength{\tabcolsep}{5pt}
\resizebox{\columnwidth}{!}{%
\begin{tabular}{ll}
\toprule
Hyperparameter & Value \\
\midrule
Input length / sample rate & 4~s (64,000 samples) / 16~kHz \\
SSL encoder & WavLM-Large, frozen, $L{=}25$, $D{=}1024$ \\
Common channel dim.\ $C$ / embedding $d$ & 128 / 160 \\
Optimizer & AdamW~\cite{loshchilov2019adamw}, wd $10^{-4}$ \\
LR (new modules / schedule) & $10^{-4}$, cosine, 2-epoch warmup \\
Epochs / batch size & 40 / 24 \\
OC-Softmax $(s_0, m_0, m_1)$ & $(20, 0.9, 0.2)$ \\
SupCon $(\lambda_c, \tau_c)$ & $(0.5, 0.1)$ \\
Entropy reg.\ $\lambda_r$ & 0.01 \\
Augmentation $p_{\text{aug}}$ / curriculum & 0.7 / 10-epoch anneal \\
Seeds & 3 (report mean $\pm$ 95\% CI) \\
\bottomrule
\end{tabular}%
}
\end{table}

Table~\ref{tab:hyper} lists hyperparameters. All audio is resampled to 16~kHz mono and
cropped or padded to 4~s. The WavLM-Large front-end is frozen; layer-attention weights,
the View-B CNN, fusion, back-end, and heads receive gradients ($\approx$4.6M
parameters). Gradient checkpointing of the frozen encoder keeps peak memory below
19~GB, enabling training on a single 24~GB GPU (RTX~3090); one run takes
$\approx$31~hours. Every experiment is repeated with three random seeds and we report
means with bootstrap confidence intervals. Baselines B1--B5 are retrained with
identical data, augmentation, seeds, and selection criteria to guarantee a controlled
comparison; B6 uses the authors' released recipe under our protocol; B7 numbers are
quoted from the original publication; B8 requires no training.

\subsection{Evaluation Metrics}
\label{ssec:metrics}

We report the metric bundle recommended by the ASVspoof community plus
calibration measures: equal error rate (EER, primary), minimum tandem detection cost
function (min~t-DCF)~\cite{kinnunen2020tdcf} with the official ASVspoof~2019 ASV
scores and cost parameters ($C_{\text{miss}}{=}1$, $C_{\text{fa}}{=}10$,
$\pi_{\text{tar}}{=}0.9405$), area under the ROC curve (AUROC), balanced accuracy,
and F1. Calibration is measured by expected calibration error (ECE, 15 equal-mass
bins)~\cite{guo2017calibration} and the Brier score, with reliability diagrams for
qualitative inspection. Statistical uncertainty is quantified with 95\% bootstrap
confidence intervals over 1,000 utterance-level resamples; system comparisons use
paired bootstrap tests on the EER difference with significance at $p < 0.05$.
Efficiency is reported as trainable/total parameters, GMACs at 4~s input, and
real-time factor on GPU and CPU (Table~\ref{tab:budget}).
